\chapter{相关工作}
\label{chap:related}

\section{集合通信运行时}

NCCL~\cite{nccl} 和 Gloo~\cite{gloo} 使用 ring、tree 和分层算法实现常用集合
通信。MSCCL~\cite{msccl}、HiCCL~\cite{hiccl} 和 TACCL~\cite{taccl} 进一步
允许编译器或用户组合发送、接收、规约和拷贝动作。这些系统决定数据如何分块及
依赖如何调度，底层传输仍接收点对点操作。\sysname 不替换 collective schedule；
它为 schedule 提供多播-聚合流，使同一 packet offset 的输入在网络中保持共同的
路径、聚合和恢复边界。

\section{在网集合通信}

SwitchML~\cite{switchml} 在可编程交换机中执行规约，并按任务划分聚合状态；
NetReduce~\cite{netreduce} 在 NIC-FPGA 的 RDMA 路径中聚合；SHARP~\cite{sharp}
由控制面建立硬件 reduction tree。这些系统以任务级资源映射隔离并发聚合。
ATP~\cite{atp} 改用共享 slot 和数据面哈希，并在冲突时由参数服务器补全；
A2TP~\cite{a2tp} 与 PA-ATP~\cite{paatp} 分别改进公平性和进度调度。AIRP~
\cite{airp} 则根据全局视图规划网络中的聚合资源。

EPIC~\cite{epic} 把交换机暴露为传输端点，可以在交换机终结连接，也可以让端点
发起的传输状态经交换机翻译到后续链路。\Cref{ssec:inc-limits} 按组外端点、连接
终结和连接翻译归纳这些组织方式。它们分别解决资源分配或连接推进，但 packet
路径、聚合位置和恢复状态仍由不同连接或控制过程维护。\sysname 用一份组播 credit
绑定 packet offset、subchannel 和逐级 slot。responder 用整组完成事件更新拥塞
控制，并在超时或 owner mismatch 时启动整组 repair。

\section{Receiver-driven 传输与拥塞控制}

Homa~\cite{homa}、NDP~\cite{ndp} 和 ExpressPass~\cite{expresspass} 由接收端
发放 grant 或 credit，以限制发送端注入。\sysname 同样把发送许可放在接收端，
但一份 credit 同时授权 requester group 对一个 packet offset 各发送一份输入。
该许可还指定 subchannel，并在返回路径取得聚合位置；普通 unicast grant 不需要
维持这些组内一致性条件。

DCTCP~\cite{dctcp}、DCQCN~\cite{dcqcn}、HPCC~\cite{hpcc} 和 TIMELY~
\cite{timely} 分别使用 ECN、INT 或 RTT 调节发送窗口和速率。\sysname 不固定
控制律，而是在 responder 提供 rate/window 两道 credit gate。基于时延的算法
读取整组 fan-in 的完成时间；ECN 算法读取各分支 CE 的 OR。该接口改变反馈的
聚合范围，但保留所选控制律自身的收敛和公平性限制（\cref{ssec:cc}）。

\section{多路径传输}

ECMP~\cite{ecmp} 按 flow 哈希路径；CONGA~\cite{conga}、Hermes~\cite{hermes}、
Presto~\cite{presto} 和 LetFlow~\cite{letflow} 根据网络状态在 flowlet 或 packet
粒度切换路径。Themis~\cite{themis} 使用确定性的 PSN--path 映射，在 commodity
RNIC 的重传语义下支持 packet spraying，并用基于乱序程度的信号暂时避开检测到
的慢路径或失效路径。MRC~\cite{mrc,mrcsrv6} 定义 per-connection EV 集合和
sender-side congestion-control 信号，并把具体 load-balancing algorithm 留给
实现；其生产实现逐 packet 轮转 active EV，临时避开带 ECN 的 EV，并在路径丢包
后替换 EV。这些机制以单个 unicast packet 或 flow 为调度对象。

\sysname 为每棵预装 subchannel 维护独立的 rate/window gate，并把后续 offset
分配给当前允许发放的候选。该选择器可以调度单个 packet；用于在网 fan-in 时，
一次决定覆盖同一 packet offset 的整组输入。整组完成事件更新所选 subchannel
的 gate，使后续 offset 只在当前允许发放的聚合树间分配。

在网 fan-in 另有一项正确性条件：同一 packet offset 的全部输入必须选择同一
聚合树，并在共享子树中命中同一 slot。\sysname 由 responder 选择 subchannel，
再通过组播 credit 把选择复制给全部 requester；已发出的 Credit Context 不迁移。
